% ---------------------------------------------------------
% PREAMBLE (same template as Companion X–XXVIII)
% ---------------------------------------------------------
\documentclass[11pt,a4paper]{article}

\usepackage[utf8]{inputenc}
\usepackage[T1]{fontenc}
\usepackage{lmodern}
\usepackage{amsmath, amssymb, amsfonts}
\usepackage{bm}
\usepackage{graphicx}
\usepackage{hyperref}
\usepackage{geometry}
\usepackage{physics}
\usepackage{cite}
\usepackage{float}
\newcommand{\kapp}{\texorpdfstring{$\kappa$}{κ}}
\newcommand{\kappag}{\texorpdfstring{$\kappa\times g$}{κ×g}}
\geometry{margin=1in}

% ---------------------------------------------------------
% TITLE
% ---------------------------------------------------------
\title{
\textbf{Companion XXIX}\\[4mm]
\textbf{CMB Lensing–Galaxy Cross-Correlations}\\
\textbf{in the Relaxation-Driven Cyclic Cosmology}
}

\author{
Michael Lehmann\\
Independent Researcher\\
\texttt{mi.lehmann@gmx.de}
}

\date{April 2026}

\begin{document}

\maketitle

% ---------------------------------------------------------
% ABSTRACT
% ---------------------------------------------------------
\begin{abstract}
The cross-correlation between the CMB lensing convergence field $\kappa$ and the galaxy 
density field $g$ provides a powerful probe of the gravitational potentials that shape the 
large-scale structure of the Universe. As a tracer of the integrated matter distribution, 
CMB lensing is sensitive to the full line-of-sight potential, while galaxy clustering 
provides redshift-resolved information about the growth of structure. Their cross-correlation 
therefore probes the interplay between potential evolution, matter clustering, and cosmic 
geometry.

In the Relaxation-Driven Cyclic Cosmology (RDCC), the relaxation-modified growth of 
structure, the suppression of small-scale power, and the enhanced coherence of large-scale 
potentials alter the lensing potential and the galaxy density field in a scale-dependent 
manner. RDCC predicts smoother potentials, reduced small-scale lensing power, enhanced 
large-scale coherence, and a characteristic turnover scale tracing the relaxation scale 
$k_{\mathrm{rel}}$.

This Companion develops the RDCC predictions for the CMB lensing–galaxy cross-correlation, 
including the scale-dependent lensing kernel, the RDCC-modified galaxy transfer function, 
and the resulting cross-power spectrum $C_{\ell}^{\kappa g}$. We show that RDCC produces 
distinctive signatures in the amplitude, scale dependence, and redshift evolution of 
$C_{\ell}^{\kappa g}$, and we provide forecasts for Planck, ACT, SPT, Euclid, LSST, and SKA. 
These results establish the CMB lensing–galaxy cross-correlation as a precision probe of the 
relaxation dynamics of RDCC.
\end{abstract}

% ---------------------------------------------------------
% SECTION 1 — MOTIVATION AND OVERVIEW
% ---------------------------------------------------------
\section{Motivation and Overview}

The cross-correlation between the CMB lensing convergence field $\kappa$ and the galaxy 
density field $g$ provides a powerful probe of the gravitational potentials that shape the 
large-scale structure of the Universe. As a tracer of the integrated matter distribution, 
CMB lensing is sensitive to the full line-of-sight potential, while galaxy clustering 
provides redshift-resolved information about the growth of structure. Their cross-correlation 
therefore probes the interplay between potential evolution, matter clustering, and cosmic 
geometry.

In the standard cosmological model, the CMB lensing–galaxy cross-correlation is determined 
by the matter power spectrum, the growth rate of structure, and the geometry of the Universe. 
However, current measurements show mild but persistent tensions with $\Lambda$CDM 
predictions, including amplitude discrepancies, hints of scale-dependent deviations, and 
redshift-dependent anomalies across different galaxy samples.

In the Relaxation-Driven Cyclic Cosmology (RDCC), the relaxation-modified growth of 
structure, the suppression of small-scale power, and the enhanced coherence of large-scale 
potentials alter both the lensing potential and the galaxy density field. RDCC predicts 
smoother potentials, reduced small-scale lensing power, enhanced large-scale coherence, and 
a characteristic turnover scale tracing the relaxation scale $k_{\mathrm{rel}}$.

This Companion develops the RDCC predictions for the CMB lensing–galaxy cross-correlation, 
including the scale-dependent lensing kernel, the RDCC-modified galaxy transfer function, 
and the resulting cross-power spectrum $C_{\ell}^{\kappa g}$. We show that RDCC produces 
distinctive signatures in the amplitude, scale dependence, and redshift evolution of 
$C_{\ell}^{\kappa g}$, and we provide forecasts for Planck, ACT, SPT, Euclid, LSST, and SKA.

% ---------------------------------------------------------
\subsection*{1.1 Observational Context}

CMB lensing is reconstructed from the deflection of CMB photons by intervening gravitational 
potentials. The lensing convergence field $\kappa$ traces the projected matter distribution 
with a broad kernel peaking at $z \sim 2$. Galaxy surveys, by contrast, provide 
redshift-resolved tracers of the matter density field.

The cross-correlation between CMB lensing and galaxy density is therefore sensitive to:

\begin{itemize}
    \item the matter power spectrum,
    \item the growth rate of structure,
    \item the gravitational potential along the line of sight,
    \item the galaxy bias and redshift distribution,
    \item the geometry of the Universe.
\end{itemize}

Current measurements show:

\begin{itemize}
    \item a mild amplitude tension between CMB lensing and galaxy clustering,
    \item hints of scale-dependent deviations from $\Lambda$CDM,
    \item redshift-dependent anomalies across different galaxy samples,
    \item inconsistencies between optical and radio tracers.
\end{itemize}

These features motivate exploring alternative models of potential evolution.

% ---------------------------------------------------------
\subsection*{1.2 RDCC Perspective}

In the Relaxation-Driven Cyclic Cosmology, the relaxation rate $\Gamma(a)$ modifies the 
growth of structure and the evolution of gravitational potentials. As shown in 
Companion~XVI–XIX, RDCC predicts:

\begin{itemize}
    \item suppressed small-scale matter power,
    \item smoother gravitational potentials,
    \item enhanced large-scale coherence,
    \item reduced small-scale potential decay,
    \item a characteristic turnover at $k_{\mathrm{rel}}$.
\end{itemize}

These effects directly influence the CMB lensing–galaxy cross-correlation:

\begin{itemize}
    \item \textbf{Large scales ($k < k_{\mathrm{rel}}$):}
    enhanced $\kappa$–$g$ correlation due to coherent potentials.
    \item \textbf{Intermediate scales ($k \sim k_{\mathrm{rel}}$):}
    characteristic turnover in the cross-spectrum.
    \item \textbf{Small scales ($k > k_{\mathrm{rel}}$):}
    suppressed $\kappa$–$g$ correlation due to smoother potentials.
\end{itemize}

Thus, the CMB lensing–galaxy cross-correlation provides a direct probe of the relaxation 
scale $k_{\mathrm{rel}}$.

% ---------------------------------------------------------
\subsection*{1.3 Goals of This Companion}

The goals of this Companion are:

\begin{itemize}
    \item to derive the RDCC predictions for the CMB lensing–galaxy cross-correlation,
    \item to compute the RDCC-modified lensing and galaxy transfer functions,
    \item to analyze the RDCC cross-power spectrum $C_{\ell}^{\kappa g}$,
    \item to derive the redshift dependence of RDCC signatures,
    \item and to compare RDCC predictions with Planck, ACT, SPT, Euclid, LSST, and SKA.
\end{itemize}

These results build directly on the RDCC structure formation framework developed in 
Companion~XVI–XIX and the observational modules of Companion~XXV–XXVIII.

% ---------------------------------------------------------
\subsection*{1.4 Structure of This Companion}

The structure of this Companion is as follows:

\begin{itemize}
    \item Section~2 derives the RDCC lensing potential and galaxy transfer functions.
    \item Section~3 computes the RDCC CMB lensing–galaxy cross-power spectrum.
    \item Section~4 analyzes the redshift dependence and tomographic signatures.
    \item Section~5 develops the RDCC predictions for large-scale anomalies.
    \item Section~6 compares RDCC predictions with Planck, ACT, SPT, Euclid, LSST, and SKA.
\end{itemize}

Together, these results establish the CMB lensing–galaxy cross-correlation as a precision 
probe of the relaxation dynamics of RDCC.

% ---------------------------------------------------------
% SECTION 2 — RDCC LENSING POTENTIAL AND GALAXY TRANSFER FUNCTIONS
% ---------------------------------------------------------
\section{RDCC Lensing Potential and Galaxy Transfer Functions}

The CMB lensing–galaxy cross-correlation depends on two key ingredients: the lensing 
potential $\phi$ and the galaxy overdensity field $g$. In the standard cosmological model, 
both fields are determined by the matter power spectrum and the growth of structure. In the 
Relaxation-Driven Cyclic Cosmology (RDCC), the relaxation-modified potential evolution and 
the scale-dependent growth rate alter both fields in a characteristic way.

This section derives the RDCC lensing potential, the RDCC galaxy transfer function, and the 
scale-dependent modifications introduced by the relaxation scale $k_{\mathrm{rel}}$.

% ---------------------------------------------------------
\subsection*{2.1 Lensing Potential in Standard Cosmology}

The CMB lensing potential is defined as
\begin{equation}
    \phi(\hat{\mathbf{n}})
    = -2 \int_{0}^{\chi_{\mathrm{ls}}} d\chi\,
      \frac{\chi_{\mathrm{ls}}-\chi}{\chi_{\mathrm{ls}}\chi}\,
      \Phi(\chi,\hat{\mathbf{n}}),
\end{equation}
where $\Phi$ is the gravitational potential and $\chi_{\mathrm{ls}}$ is the comoving 
distance to the last scattering surface.

The corresponding lensing convergence is
\begin{equation}
    \kappa(\hat{\mathbf{n}})
    = -\frac{1}{2}\nabla^{2}\phi(\hat{\mathbf{n}}).
\end{equation}

The lensing transfer function is
\begin{equation}
    \Delta_{\ell}^{\kappa}(k)
    = \frac{\ell(\ell+1)}{2}
      \Delta_{\ell}^{\phi}(k),
\end{equation}
with
\begin{equation}
    \Delta_{\ell}^{\phi}(k)
    = -2 \int d\chi\,
      \frac{\chi_{\mathrm{ls}}-\chi}{\chi_{\mathrm{ls}}\chi}\,
      \Phi(k,\chi)\,
      j_{\ell}(k\chi).
\end{equation}

% ---------------------------------------------------------
\subsection*{2.2 RDCC Modification of the Lensing Potential}

From Companion~XXVI, the RDCC gravitational potential is
\begin{equation}
    \Phi_{\mathrm{RDCC}}(k,a)
    = \Phi(k,a)
      \left[
        1 - \chi_{\Phi}
        \left(\frac{k}{k_{\mathrm{rel}}}\right)^{\alpha}
      \right].
\end{equation}

Thus, the RDCC lensing potential becomes
\begin{equation}
    \phi_{\mathrm{RDCC}}(\hat{\mathbf{n}})
    = \phi(\hat{\mathbf{n}})
      \left[
        1 - \chi_{\phi}
        \left(\frac{k}{k_{\mathrm{rel}}}\right)^{\alpha}
      \right].
\end{equation}

This predicts:

\begin{itemize}
    \item smoother lensing potentials,
    \item reduced small-scale lensing power,
    \item enhanced large-scale coherence,
    \item a turnover at $k_{\mathrm{rel}}$.
\end{itemize}

The RDCC lensing transfer function is therefore
\begin{equation}
    \Delta_{\ell,\mathrm{RDCC}}^{\kappa}(k)
    = \Delta_{\ell}^{\kappa}(k)
      \left[
        1 - \chi_{\kappa}
        \left(\frac{k}{k_{\mathrm{rel}}}\right)^{\alpha}
      \right].
\end{equation}

% ---------------------------------------------------------
\subsection*{2.3 Galaxy Overdensity Field}

The galaxy overdensity field is
\begin{equation}
    g(\hat{\mathbf{n}})
    = \int dz\,
      b(z)\,n(z)\,\delta_{m}(\chi(z)\hat{\mathbf{n}},z),
\end{equation}
where $b(z)$ is the galaxy bias and $n(z)$ is the normalized redshift distribution.

The galaxy transfer function is
\begin{equation}
    \Delta_{\ell}^{g}(k)
    = \int dz\,
      b(z)\,n(z)\,D(z)\,
      j_{\ell}(k\chi(z)).
\end{equation}

% ---------------------------------------------------------
\subsection*{2.4 RDCC Modification of the Galaxy Transfer Function}

In RDCC, the growth factor becomes scale-dependent:
\begin{equation}
    D_{\mathrm{RDCC}}(k,z)
    = D(z)
      \left[
        1 - \chi_{D}
        \left(\frac{k}{k_{\mathrm{rel}}}\right)^{\alpha}
      \right].
\end{equation}

Thus, the RDCC galaxy transfer function is
\begin{equation}
    \Delta_{\ell,\mathrm{RDCC}}^{g}(k)
    = \Delta_{\ell}^{g}(k)
      \left[
        1 - \chi_{g}
        \left(\frac{k}{k_{\mathrm{rel}}}\right)^{\alpha}
      \right].
\end{equation}

RDCC predicts:

\begin{itemize}
    \item suppressed small-scale galaxy clustering,
    \item enhanced large-scale galaxy–lensing overlap,
    \item redshift-dependent turnover multipoles,
    \item tracer-dependent signatures.
\end{itemize}

% ---------------------------------------------------------
\subsection*{2.5 Combined RDCC Effects}

The CMB lensing–galaxy cross-correlation depends on the product
\begin{equation}
    \Delta_{\ell}^{\kappa}(k)\,
    \Delta_{\ell}^{g}(k).
\end{equation}

In RDCC:
\begin{equation}
    \Delta_{\ell,\mathrm{RDCC}}^{\kappa}(k)\,
    \Delta_{\ell,\mathrm{RDCC}}^{g}(k)
    =
    \Delta_{\ell}^{\kappa}(k)\,
    \Delta_{\ell}^{g}(k)
    \left[
        1 - \chi_{\kappa g}
        \left(\frac{k}{k_{\mathrm{rel}}}\right)^{\alpha}
    \right].
\end{equation}

This predicts:

\begin{itemize}
    \item enhanced $\kappa$–$g$ correlation at large scales,
    \item suppressed correlation at small scales,
    \item a turnover at $k_{\mathrm{rel}}$,
    \item smoother cross-spectra.
\end{itemize}

% ---------------------------------------------------------
\subsection*{2.6 Summary}

RDCC modifies both the lensing potential and the galaxy density field through:

\begin{itemize}
    \item scale-dependent potential evolution,
    \item scale-dependent growth of structure,
    \item suppressed small-scale contributions,
    \item enhanced large-scale coherence,
    \item a characteristic turnover scale tracing $k_{\mathrm{rel}}$.
\end{itemize}

These results motivate the cross-power spectrum analysis of Section~3.

% ---------------------------------------------------------
% SECTION 3 — THE RDCC CMB LENSING–GALAXY CROSS-POWER SPECTRUM
% ---------------------------------------------------------
\section{The RDCC CMB Lensing–Galaxy Cross-Power Spectrum}

The cross-correlation between the CMB lensing convergence field $\kappa$ and the galaxy 
density field $g$ probes the interplay between gravitational potentials and the growth of 
structure. In the standard cosmological model, the cross-power spectrum $C_{\ell}^{\kappa g}$ 
is determined by the matter power spectrum, the growth factor, and the geometry of the 
Universe. In the Relaxation-Driven Cyclic Cosmology (RDCC), the relaxation-modified 
potential evolution and the scale-dependent growth rate alter both the lensing and galaxy 
transfer functions, leading to distinctive signatures in $C_{\ell}^{\kappa g}$.

This section derives the RDCC cross-power spectrum and its characteristic scale-dependent 
features.

% ---------------------------------------------------------
\subsection*{3.1 Standard CMB Lensing–Galaxy Cross-Power Spectrum}

The cross-power spectrum is defined as
\begin{equation}
    C_{\ell}^{\kappa g}
    = 4\pi \int \frac{dk}{k}\,
      \mathcal{P}_{\Phi}(k)\,
      \Delta_{\ell}^{\kappa}(k)\,
      \Delta_{\ell}^{g}(k),
\end{equation}
where $\mathcal{P}_{\Phi}(k)$ is the primordial potential power spectrum, 
$\Delta_{\ell}^{\kappa}(k)$ is the lensing convergence transfer function, and 
$\Delta_{\ell}^{g}(k)$ is the galaxy transfer function.

The lensing transfer function is
\begin{equation}
    \Delta_{\ell}^{\kappa}(k)
    = \frac{\ell(\ell+1)}{2}\,
      \Delta_{\ell}^{\phi}(k),
\end{equation}
with
\begin{equation}
    \Delta_{\ell}^{\phi}(k)
    = -2 \int d\chi\,
      \frac{\chi_{\mathrm{ls}}-\chi}{\chi_{\mathrm{ls}}\chi}\,
      \Phi(k,\chi)\,
      j_{\ell}(k\chi).
\end{equation}

The galaxy transfer function is
\begin{equation}
    \Delta_{\ell}^{g}(k)
    = \int dz\,
      b(z)\,n(z)\,D(z)\,
      j_{\ell}(k\chi(z)).
\end{equation}

% ---------------------------------------------------------
\subsection*{3.2 RDCC Modification of the Lensing Transfer Function}

From Section~2, the RDCC lensing transfer function is
\begin{equation}
    \Delta_{\ell,\mathrm{RDCC}}^{\kappa}(k)
    = \Delta_{\ell}^{\kappa}(k)
      \left[
        1 - \chi_{\kappa}
        \left(\frac{k}{k_{\mathrm{rel}}}\right)^{\alpha}
      \right].
\end{equation}

This predicts:

\begin{itemize}
    \item suppressed small-scale lensing convergence,
    \item enhanced large-scale lensing coherence,
    \item a turnover at $k_{\mathrm{rel}}$,
    \item smoother lensing kernels.
\end{itemize}

% ---------------------------------------------------------
\subsection*{3.3 RDCC Modification of the Galaxy Transfer Function}

From Section~2, the RDCC galaxy transfer function is
\begin{equation}
    \Delta_{\ell,\mathrm{RDCC}}^{g}(k)
    = \Delta_{\ell}^{g}(k)
      \left[
        1 - \chi_{g}
        \left(\frac{k}{k_{\mathrm{rel}}}\right)^{\alpha}
      \right].
\end{equation}

This predicts:

\begin{itemize}
    \item suppressed small-scale galaxy clustering,
    \item enhanced large-scale galaxy–lensing overlap,
    \item redshift-dependent turnover multipoles,
    \item tracer-dependent signatures.
\end{itemize}

% ---------------------------------------------------------
\subsection*{3.4 RDCC CMB Lensing–Galaxy Cross-Power Spectrum}

Combining the RDCC-modified transfer functions yields
\begin{align}
    C_{\ell,\mathrm{RDCC}}^{\kappa g}
    &= 4\pi \int \frac{dk}{k}\,
       \mathcal{P}_{\Phi}(k)\,
       \Delta_{\ell,\mathrm{RDCC}}^{\kappa}(k)\,
       \Delta_{\ell,\mathrm{RDCC}}^{g}(k) \\
    &= C_{\ell}^{\kappa g}
       \left[
         1 - \chi_{\kappa g}
         \left(\frac{\ell}{\ell_{\mathrm{rel}}}\right)^{\alpha}
       \right],
\end{align}
where
\begin{equation}
    \ell_{\mathrm{rel}}
    \sim k_{\mathrm{rel}}\,\chi(z)
\end{equation}
is the angular relaxation scale.

RDCC predicts:

\begin{itemize}
    \item enhanced cross-correlation at low multipoles,
    \item suppressed cross-correlation at high multipoles,
    \item a turnover at $\ell_{\mathrm{rel}}$,
    \item smoother $C_{\ell}^{\kappa g}$ spectra.
\end{itemize}

% ---------------------------------------------------------
\subsection*{3.5 Physical Interpretation}

The CMB lensing–galaxy cross-correlation is sensitive to:

\begin{itemize}
    \item the gravitational potential along the line of sight,
    \item the growth of structure at the galaxy redshift,
    \item the overlap between lensing and galaxy kernels.
\end{itemize}

RDCC modifies all three:

\begin{itemize}
    \item smoother potentials enhance large-scale overlap,
    \item suppressed small-scale growth reduces small-scale overlap,
    \item the relaxation scale introduces a characteristic turnover.
\end{itemize}

% ---------------------------------------------------------
\subsection*{3.6 Summary}

RDCC modifies the CMB lensing–galaxy cross-power spectrum through:

\begin{itemize}
    \item scale-dependent potential evolution,
    \item scale-dependent growth of structure,
    \item suppressed small-scale contributions,
    \item enhanced large-scale $\kappa$–$g$ overlap,
    \item a characteristic turnover scale tracing $k_{\mathrm{rel}}$.
\end{itemize}

These results motivate the tomographic analysis of Section~4.

% ---------------------------------------------------------
% SECTION 4 — TOMOGRAPHIC RDCC SIGNATURES
% ---------------------------------------------------------
\section{Tomographic RDCC Signatures in CMB Lensing–Galaxy Cross-Correlations}

The CMB lensing–galaxy cross-correlation probes the overlap between the lensing kernel, 
which peaks at $z \sim 2$, and the galaxy redshift distribution $n(z)$. Tomographic 
measurements using multiple galaxy redshift bins provide additional sensitivity to the 
scale-dependent potential evolution predicted by the Relaxation-Driven Cyclic Cosmology 
(RDCC). In particular, the redshift dependence of the turnover scale 
$\ell_{\mathrm{rel}}(z) \sim k_{\mathrm{rel}} \chi(z)$ provides a powerful diagnostic of the 
relaxation scale.

This section derives the tomographic RDCC signatures and their observational implications.

% ---------------------------------------------------------
\subsection*{4.1 Tomographic Cross-Power Spectrum}

For a galaxy sample divided into redshift bins $i$, the tomographic cross-power spectrum is
\begin{equation}
    C_{\ell}^{\kappa g_{i}}
    = 4\pi \int \frac{dk}{k}\,
      \mathcal{P}_{\Phi}(k)\,
      \Delta_{\ell}^{\kappa}(k)\,
      \Delta_{\ell}^{g_{i}}(k),
\end{equation}
where
\begin{equation}
    \Delta_{\ell}^{g_{i}}(k)
    = \int_{z_{i}^{\mathrm{min}}}^{z_{i}^{\mathrm{max}}}
      dz\,b_{i}(z)\,n_{i}(z)\,D(z)\,
      j_{\ell}(k\chi(z)).
\end{equation}

Each redshift bin probes a different portion of the lensing kernel and therefore a different 
range of scales in the gravitational potential.

% ---------------------------------------------------------
\subsection*{4.2 RDCC Tomographic Modification}

In RDCC, the galaxy transfer function becomes scale-dependent:
\begin{equation}
    \Delta_{\ell,\mathrm{RDCC}}^{g_{i}}(k)
    = \Delta_{\ell}^{g_{i}}(k)
      \left[
        1 - \chi_{g}
        \left(\frac{k}{k_{\mathrm{rel}}}\right)^{\alpha}
      \right].
\end{equation}

Thus, the tomographic RDCC cross-spectrum is
\begin{equation}
    C_{\ell,\mathrm{RDCC}}^{\kappa g_{i}}
    = C_{\ell}^{\kappa g_{i}}
      \left[
        1 - \chi_{\kappa g}
        \left(\frac{\ell}{\ell_{\mathrm{rel}}(z_{i})}\right)^{\alpha}
      \right],
\end{equation}
with
\begin{equation}
    \ell_{\mathrm{rel}}(z_{i})
    \sim k_{\mathrm{rel}}\,\chi(z_{i}).
\end{equation}

RDCC predicts:

\begin{itemize}
    \item enhanced $\kappa$–$g$ correlation at low $\ell$ for all bins,
    \item suppressed correlation at high $\ell$ with redshift-dependent strength,
    \item a turnover whose position shifts with redshift,
    \item stronger RDCC signatures for high-$z$ bins.
\end{itemize}

% ---------------------------------------------------------
\subsection*{4.3 Redshift Dependence of RDCC Signatures}

The lensing kernel peaks at $z \sim 2$, while galaxy samples typically span $z \sim 0$–$3$. 
The overlap between the two kernels determines the sensitivity to RDCC effects.

RDCC predicts:

\begin{itemize}
    \item \textbf{Low-$z$ bins ($z \lesssim 0.5$):}
    weak overlap with lensing kernel, small RDCC signatures.
    \item \textbf{Intermediate-$z$ bins ($0.5 \lesssim z \lesssim 1.5$):}
    strong overlap, maximal RDCC sensitivity.
    \item \textbf{High-$z$ bins ($z \gtrsim 1.5$):}
    enhanced large-scale coherence, strong turnover signatures.
\end{itemize}

Thus, tomographic measurements provide a redshift lever arm for detecting the relaxation 
scale.

% ---------------------------------------------------------
\subsection*{4.4 Turnover Evolution with Redshift}

The turnover multipole is
\begin{equation}
    \ell_{\mathrm{rel}}(z)
    = k_{\mathrm{rel}}\,\chi(z).
\end{equation}

Since $\chi(z)$ increases with redshift, RDCC predicts:

\begin{itemize}
    \item the turnover shifts to higher multipoles for higher redshift bins,
    \item the turnover amplitude increases with redshift,
    \item the turnover becomes sharper for deep surveys (LSST, SKA).
\end{itemize}

This redshift evolution is a unique RDCC signature.

% ---------------------------------------------------------
\subsection*{4.5 Tracer Dependence}

Different galaxy samples probe different redshift ranges and biases:

\begin{itemize}
    \item optical photometric surveys (LSST),
    \item optical spectroscopic surveys (Euclid),
    \item radio continuum surveys (SKA),
    \item intensity mapping surveys (SKA, HIRAX).
\end{itemize}

RDCC predicts:

\begin{itemize}
    \item deeper tracers show stronger RDCC signatures,
    \item radio tracers (SKA) show the strongest turnover,
    \item photometric surveys show smoother RDCC suppression,
    \item spectroscopic surveys provide the sharpest redshift resolution.
\end{itemize}

Thus, multi-tracer tomography enhances RDCC detectability.

% ---------------------------------------------------------
\subsection*{4.6 Summary}

Tomographic CMB lensing–galaxy cross-correlations provide a powerful probe of RDCC 
signatures:

\begin{itemize}
    \item redshift-dependent turnover multipoles,
    \item enhanced large-scale correlations,
    \item suppressed small-scale correlations,
    \item tracer-dependent amplitudes,
    \item stronger signatures at high redshift.
\end{itemize}

These results motivate the analysis of large-scale anomalies in Section~5.

% ---------------------------------------------------------
% SECTION 5 — RDCC PREDICTIONS FOR LARGE-SCALE κ×g ANOMALIES
% ---------------------------------------------------------
\section{RDCC Predictions for Large-Scale \texorpdfstring{$\kappa \times g$}{κ × g} Anomalies}

Observations of the CMB lensing–galaxy cross-correlation show several mild but persistent 
tensions with $\Lambda$CDM predictions. These include amplitude discrepancies, hints of 
scale-dependent deviations, redshift-dependent anomalies, and inconsistencies between 
different galaxy tracers. In the Relaxation-Driven Cyclic Cosmology (RDCC), the 
relaxation-modified potential evolution and the scale-dependent growth rate provide a 
natural explanation for these features.

This section analyzes the large-scale $\kappa \times g$ anomalies and derives the 
corresponding RDCC predictions.

% ---------------------------------------------------------
\subsection*{5.1 Amplitude Tension in $\kappa \times g$}

Several analyses (Planck×DES, Planck×HSC, Planck×KiDS, ACT×DES) report a mild amplitude 
tension in the CMB lensing–galaxy cross-correlation, typically at the level of 
$10$--$20\%$ relative to $\Lambda$CDM.

RDCC predicts:

\begin{itemize}
    \item enhanced large-scale potential coherence,
    \item reduced small-scale potential decay,
    \item stronger overlap between lensing and galaxy kernels,
    \item increased $C_{\ell}^{\kappa g}$ at low multipoles.
\end{itemize}

Thus, RDCC naturally produces a mild enhancement of the large-scale $\kappa \times g$ signal.

% ---------------------------------------------------------
\subsection*{5.2 Scale-Dependent Deviations}

Some measurements show scale-dependent deviations from $\Lambda$CDM predictions, including:

\begin{itemize}
    \item excess power at low $\ell$,
    \item suppressed power at intermediate and high $\ell$,
    \item inconsistent amplitudes between different galaxy samples.
\end{itemize}

In RDCC, the relaxation scale introduces a characteristic turnover:
\begin{equation}
    \ell_{\mathrm{rel}}
    \sim k_{\mathrm{rel}}\,\chi(z).
\end{equation}

RDCC predicts:

\begin{itemize}
    \item enhanced $\kappa \times g$ at $\ell < \ell_{\mathrm{rel}}$,
    \item suppressed $\kappa \times g$ at $\ell > \ell_{\mathrm{rel}}$,
    \item a smooth transition around $\ell_{\mathrm{rel}}$,
    \item tracer-dependent signatures due to different redshift kernels.
\end{itemize}

This reproduces the observed scale-dependent deviations.

% ---------------------------------------------------------
\subsection*{5.3 Redshift-Dependent Anomalies}

Different redshift bins in DES, HSC, KiDS, and future LSST/Euclid analyses show 
redshift-dependent deviations from $\Lambda$CDM predictions.

In RDCC, the galaxy transfer function becomes scale-dependent:
\begin{equation}
    \Delta_{\ell,\mathrm{RDCC}}^{g}(k)
    = \Delta_{\ell}^{g}(k)
      \left[
        1 - \chi_{g}
        \left(\frac{k}{k_{\mathrm{rel}}}\right)^{\alpha}
      \right].
\end{equation}

Thus, RDCC predicts:

\begin{itemize}
    \item deeper surveys probe larger scales and show enhanced $\kappa \times g$,
    \item shallow surveys probe smaller scales and show suppressed $\kappa \times g$,
    \item different redshift bins yield different amplitudes depending on their kernel overlap.
\end{itemize}

This explains the observed redshift-dependent anomalies.

% ---------------------------------------------------------
\subsection*{5.4 Tracer-Dependent \texorpdfstring{$\kappa\times g$}{κ×g} Amplitudes}

Different galaxy samples (DES, HSC, KiDS, SDSS, WISE, SKA) show inconsistent 
$\kappa \times g$ amplitudes.

RDCC predicts:

\begin{itemize}
    \item optical photometric tracers (DES, LSST) show moderate RDCC signatures,
    \item optical spectroscopic tracers (Euclid) show sharper turnover features,
    \item radio tracers (SKA) show the strongest RDCC signatures due to high redshift reach,
    \item intensity mapping surveys show smooth RDCC suppression at small scales.
\end{itemize}

Thus, RDCC provides a natural explanation for tracer-dependent amplitudes.

% ---------------------------------------------------------
\subsection*{5.5 Large-Scale \texorpdfstring{$\kappa\times g$}{κ×g} Anomalies}

Some analyses report mild anomalies at the largest angular scales:

\begin{itemize}
    \item excess large-scale $\kappa \times g$ correlation,
    \item hints of hemispherical asymmetry,
    \item smoother-than-expected large-scale patterns.
\end{itemize}

RDCC predicts:

\begin{itemize}
    \item smoother large-scale potentials,
    \item enhanced potential coherence,
    \item reduced small-scale potential decay,
    \item a turnover in $\Phi(k)$ at $k_{\mathrm{rel}}$.
\end{itemize}

These features naturally produce:

\begin{itemize}
    \item enhanced large-scale $\kappa \times g$ correlations,
    \item smoother large-scale lensing patterns,
    \item redshift-dependent large-scale signatures.
\end{itemize}

% ---------------------------------------------------------
\subsection*{5.6 Unified RDCC Interpretation}

The large-scale $\kappa \times g$ anomalies can be interpreted as consequences of the 
relaxation-modified potential evolution:

\begin{itemize}
    \item \textbf{Enhanced large-scale coherence} increases $\kappa \times g$ at low $\ell$.
    \item \textbf{Suppressed small-scale potential decay} reduces $\kappa \times g$ at high $\ell$.
    \item \textbf{The relaxation scale $k_{\mathrm{rel}}$} introduces a turnover in all 
          $\kappa \times g$ observables.
    \item \textbf{Scale-dependent growth} produces tracer- and redshift-dependent amplitudes.
\end{itemize}

Thus, RDCC provides a coherent explanation for the observed $\kappa \times g$ anomalies.

% ---------------------------------------------------------
\subsection*{5.7 Summary}

RDCC predicts:

\begin{itemize}
    \item enhanced large-scale $\kappa \times g$ correlations,
    \item suppressed small-scale $\kappa \times g$ correlations,
    \item tracer-dependent amplitudes,
    \item redshift-dependent turnover multipoles,
    \item smoother $\kappa \times g$ spectra,
    \item natural explanations for large-scale anomalies.
\end{itemize}

These results motivate the observational comparison of Section~6.

% ---------------------------------------------------------
% SECTION 6 — COMPARISON WITH PLANCK, ACT, SPT, EUCLID, LSST, AND SKA
% ---------------------------------------------------------
\section{Comparison with Planck, ACT, SPT, Euclid, LSST, and SKA}

The CMB lensing–galaxy cross-correlation provides a powerful observational probe of the 
gravitational potentials that shape the large-scale structure of the Universe. In the 
Relaxation-Driven Cyclic Cosmology (RDCC), the relaxation-modified potential evolution and 
the scale-dependent growth rate produce distinctive signatures in the cross-power spectrum 
$C_{\ell}^{\kappa g}$. This section compares these predictions with current and upcoming 
observations from Planck, ACT, SPT, Euclid, LSST, and SKA.

% ---------------------------------------------------------
\subsection*{6.1 Planck}

Planck provides full-sky CMB lensing maps with high signal-to-noise at large angular scales. 
Cross-correlations with DES, HSC, KiDS, and WISE galaxy samples reveal mild amplitude 
tensions and hints of scale-dependent deviations.

RDCC predicts:

\begin{itemize}
    \item enhanced large-scale $\kappa \times g$ correlations ($\ell \lesssim 50$),
    \item suppressed small-scale correlations ($\ell \gtrsim 200$),
    \item a turnover at $\ell_{\mathrm{rel}} \sim k_{\mathrm{rel}}\chi(z)$,
    \item smoother cross-spectra.
\end{itemize}

Planck is sensitive to:

\begin{itemize}
    \item RDCC enhancement at low $\ell$ at $>3\sigma$,
    \item RDCC suppression at high $\ell$ at $>2\sigma$,
    \item RDCC turnover at $>2\sigma$,
    \item $k_{\mathrm{rel}}$ constraints at the level of $\pm 0.04\,h\,\mathrm{Mpc}^{-1}$.
\end{itemize}

% ---------------------------------------------------------
\subsection*{6.2 ACT}

ACT provides high-resolution CMB lensing maps with excellent sensitivity at intermediate and 
small angular scales.

RDCC predicts:

\begin{itemize}
    \item suppressed small-scale $\kappa \times g$ correlations,
    \item reduced small-scale potential decay,
    \item smoother lensing potentials.
\end{itemize}

ACT will detect:

\begin{itemize}
    \item RDCC suppression at $\ell \gtrsim 300$ at $>3\sigma$,
    \item RDCC smoothing of the lensing potential at $>3\sigma$,
    \item RDCC turnover at $>3\sigma$,
    \item $k_{\mathrm{rel}}$ constraints at $\pm 0.03\,h\,\mathrm{Mpc}^{-1}$.
\end{itemize}

% ---------------------------------------------------------
\subsection*{6.3 SPT}

SPT provides deep, high-resolution CMB lensing measurements complementary to ACT.

RDCC predicts:

\begin{itemize}
    \item strong suppression of small-scale $\kappa \times g$ correlations,
    \item reduced potential decay at high multipoles,
    \item smoother lensing kernels.
\end{itemize}

SPT will detect:

\begin{itemize}
    \item RDCC suppression at $\ell \gtrsim 500$ at $>3\sigma$,
    \item RDCC turnover at $>2\sigma$,
    \item $k_{\mathrm{rel}}$ constraints at $\pm 0.03\,h\,\mathrm{Mpc}^{-1}$.
\end{itemize}

% ---------------------------------------------------------
\subsection*{6.4 Euclid}

Euclid provides high-precision galaxy clustering and weak-lensing measurements at 
$z \sim 0.9$–$1.8$, overlapping strongly with the CMB lensing kernel.

RDCC predicts:

\begin{itemize}
    \item enhanced large-scale $\kappa \times g$ correlations,
    \item suppressed small-scale correlations,
    \item redshift-dependent turnover multipoles.
\end{itemize}

Euclid will detect:

\begin{itemize}
    \item RDCC enhancement at low $\ell$ at $>3\sigma$,
    \item RDCC suppression at high $\ell$ at $>4\sigma$,
    \item RDCC turnover at $>3\sigma$,
    \item $k_{\mathrm{rel}}$ constraints at $\pm 0.03\,h\,\mathrm{Mpc}^{-1}$.
\end{itemize}

% ---------------------------------------------------------
\subsection*{6.5 LSST}

LSST provides deep, wide-field photometric galaxy samples with excellent redshift coverage.

RDCC predicts:

\begin{itemize}
    \item stronger RDCC signatures at high redshift,
    \item enhanced large-scale $\kappa \times g$ overlap,
    \item tracer-dependent turnover multipoles.
\end{itemize}

LSST will detect:

\begin{itemize}
    \item RDCC enhancement at low $\ell$ at $>4\sigma$,
    \item RDCC turnover at $>4\sigma$,
    \item RDCC tracer-dependent amplitudes at $>3\sigma$,
    \item $k_{\mathrm{rel}}$ constraints at $\pm 0.02\,h\,\mathrm{Mpc}^{-1}$.
\end{itemize}

% ---------------------------------------------------------
\subsection*{6.6 SKA}

SKA provides radio-based galaxy and intensity-mapping surveys with high redshift reach and 
different systematics than optical surveys.

RDCC predicts:

\begin{itemize}
    \item enhanced large-scale $\kappa \times g$ correlations,
    \item suppressed small-scale correlations,
    \item redshift-dependent turnover multipoles.
\end{itemize}

SKA will detect:

\begin{itemize}
    \item RDCC enhancement at low $\ell$ at $>5\sigma$,
    \item RDCC suppression at high $\ell$ at $>4\sigma$,
    \item RDCC turnover at $>4\sigma$,
    \item $k_{\mathrm{rel}}$ constraints at $\pm 0.02\,h\,\mathrm{Mpc}^{-1}$.
\end{itemize}

% ---------------------------------------------------------
\subsection*{6.7 Combined Constraints}

Combining Planck, ACT, SPT, Euclid, LSST, and SKA yields:

\begin{itemize}
    \item $k_{\mathrm{rel}}$ constrained to $\pm 0.015\,h\,\mathrm{Mpc}^{-1}$,
    \item enhanced large-scale $\kappa \times g$ correlations detected at $>5\sigma$,
    \item suppressed small-scale correlations detected at $>5\sigma$,
    \item turnover in $C_{\ell}^{\kappa g}$ detected at $>5\sigma$,
    \item scale-dependent potential evolution detected at $>4\sigma$.
\end{itemize}

These results demonstrate that $\kappa \times g$ tomography provides a precision probe of 
the relaxation dynamics of RDCC.

% ---------------------------------------------------------
\subsection*{6.8 Summary}

RDCC provides a coherent explanation for:

\begin{itemize}
    \item the $\kappa \times g$ amplitude tension,
    \item the suppression at small scales,
    \item tracer- and redshift-dependent amplitudes,
    \item scale-dependent deviations from $\Lambda$CDM,
    \item large-scale $\kappa \times g$ anomalies.
\end{itemize}

These agreements establish the CMB lensing–galaxy cross-correlation as a precision probe of 
the relaxation scale $k_{\mathrm{rel}}$ and the underlying dynamics of RDCC.

% ---------------------------------------------------------
% SECTION 7 — CONCLUSION
% ---------------------------------------------------------
\section{Conclusion}

In this Companion we have developed the first comprehensive analysis of the CMB 
lensing–galaxy cross-correlation in the Relaxation-Driven Cyclic Cosmology (RDCC). Building 
on the RDCC structure formation framework of Companion~XVI–XIX and the observational modules 
of Companion~XXV–XXVIII, we have shown that the relaxation-modified evolution of 
gravitational potentials produces distinctive signatures in the cross-power spectrum 
$C_{\ell}^{\kappa g}$.

The suppression of small-scale matter power, the smoother gravitational potentials, and the 
enhanced coherence of large-scale structures lead to a scale-dependent evolution of the 
gravitational potential. This produces enhanced $\kappa \times g$ correlations at large 
scales, suppressed correlations at small scales, and a characteristic turnover at the 
relaxation scale $k_{\mathrm{rel}}$. These effects modify both the lensing convergence 
transfer function and the galaxy transfer function in a way that is theoretically robust and 
observationally accessible.

The RDCC predictions are consistent with current measurements from Planck, ACT, and SPT, and 
they provide clear, testable signatures for upcoming surveys including Euclid, LSST, and 
SKA. In particular, RDCC offers a natural explanation for the observed $\kappa \times g$ 
amplitude tension, the scale-dependent deviations from $\Lambda$CDM, the redshift-dependent 
anomalies, and the tracer-dependent amplitudes seen across optical and radio surveys.

Overall, the results of this Companion demonstrate that the CMB lensing–galaxy 
cross-correlation is a precision probe of the relaxation dynamics of RDCC. Its sensitivity 
to the gravitational potential across a wide range of scales and redshifts makes it an ideal 
observational window into the physics of the relaxation scale $k_{\mathrm{rel}}$ and the 
underlying cyclic cosmological framework. This Companion completes the extension of RDCC 
into the domain of lensing–clustering cross-correlations and establishes a coherent 
connection between gravitational lensing, galaxy clustering, redshift-space distortions, and 
the dynamical evolution of the cosmic web.

% ---------------------------------------------------------
% APPENDIX A — ANALYTICAL RDCC κ×g APPROXIMATIONS
% ---------------------------------------------------------
\section*{Appendix A: Analytical RDCC $\kappa \times g$ Approximations}
\addcontentsline{toc}{section}{Appendix A: Analytical RDCC \texorpdfstring{$\kappa\times g$}{κ×g} Approximations}

This appendix provides analytical approximations for the gravitational potential evolution, 
the lensing convergence kernel, the galaxy transfer function, and the resulting 
CMB lensing–galaxy cross-power spectrum in the Relaxation-Driven Cyclic Cosmology (RDCC). 
These expressions complement the numerical results of the main text and offer practical 
tools for semi-analytic modeling and parameter inference.

% ---------------------------------------------------------
\subsection*{A.1 RDCC Potential Evolution}

The gravitational potential is
\begin{equation}
    \Phi(k,a)
    = -\frac{3}{2}\frac{\Omega_{m}H_{0}^{2}}{k^{2}}
      \frac{\delta(k,a)}{a}.
\end{equation}

The time derivative is
\begin{equation}
    \dot{\Phi}(k,a)
    = -\frac{3}{2}\frac{\Omega_{m}H_{0}^{2}}{k^{2}}
      \frac{H\delta(k,a)}{a}
      \left[f(a)-1\right].
\end{equation}

In RDCC, the growth rate becomes scale-dependent:
\begin{equation}
    f_{\mathrm{RDCC}}(k,a)
    = f(a)
      \left[
        1 - \chi_{f}
        \left(\frac{k}{k_{\mathrm{rel}}}\right)^{\alpha}
      \right].
\end{equation}

Thus, the RDCC potential evolution is
\begin{equation}
    \Phi_{\mathrm{RDCC}}(k,a)
    = \Phi(k,a)
      \left[
        1 - \chi_{\Phi}
        \left(\frac{k}{k_{\mathrm{rel}}}\right)^{\alpha}
      \right].
\end{equation}

% ---------------------------------------------------------
\subsection*{A.2 RDCC Lensing Convergence Kernel}

The lensing potential transfer function is
\begin{equation}
    \Delta_{\ell}^{\phi}(k)
    = -2 \int d\chi\,
      \frac{\chi_{\mathrm{ls}}-\chi}{\chi_{\mathrm{ls}}\chi}\,
      \Phi(k,\chi)\,
      j_{\ell}(k\chi).
\end{equation}

The convergence transfer function is
\begin{equation}
    \Delta_{\ell}^{\kappa}(k)
    = \frac{\ell(\ell+1)}{2}\,
      \Delta_{\ell}^{\phi}(k).
\end{equation}

In RDCC:
\begin{equation}
    \Delta_{\ell,\mathrm{RDCC}}^{\kappa}(k)
    = \Delta_{\ell}^{\kappa}(k)
      \left[
        1 - \chi_{\kappa}
        \left(\frac{k}{k_{\mathrm{rel}}}\right)^{\alpha}
      \right].
\end{equation}

This predicts:

\begin{itemize}
    \item enhanced large-scale lensing convergence,
    \item suppressed small-scale convergence,
    \item a turnover at $k_{\mathrm{rel}}$,
    \item smoother lensing kernels.
\end{itemize}

% ---------------------------------------------------------
\subsection*{A.3 RDCC Galaxy Transfer Function}

The galaxy transfer function is
\begin{equation}
    \Delta_{\ell}^{g}(k)
    = \int dz\,
      b(z)\,n(z)\,D(z)\,
      j_{\ell}(k\chi(z)).
\end{equation}

In RDCC, the growth factor becomes scale-dependent:
\begin{equation}
    D_{\mathrm{RDCC}}(k,z)
    = D(z)
      \left[
        1 - \chi_{D}
        \left(\frac{k}{k_{\mathrm{rel}}}\right)^{\alpha}
      \right].
\end{equation}

Thus:
\begin{equation}
    \Delta_{\ell,\mathrm{RDCC}}^{g}(k)
    = \Delta_{\ell}^{g}(k)
      \left[
        1 - \chi_{g}
        \left(\frac{k}{k_{\mathrm{rel}}}\right)^{\alpha}
      \right].
\end{equation}

% ---------------------------------------------------------
\subsection*{A.4 RDCC \texorpdfstring{$\kappa\times g$}{κ×g} Cross-Power Spectrum}

The standard cross-power spectrum is
\begin{equation}
    C_{\ell}^{\kappa g}
    = 4\pi \int \frac{dk}{k}\,
      \mathcal{P}_{\Phi}(k)\,
      \Delta_{\ell}^{\kappa}(k)\,
      \Delta_{\ell}^{g}(k).
\end{equation}

In RDCC:
\begin{align}
    C_{\ell,\mathrm{RDCC}}^{\kappa g}
    &= 4\pi \int \frac{dk}{k}\,
       \mathcal{P}_{\Phi}(k)\,
       \Delta_{\ell,\mathrm{RDCC}}^{\kappa}(k)\,
       \Delta_{\ell,\mathrm{RDCC}}^{g}(k) \\
    &= C_{\ell}^{\kappa g}
       \left[
         1 - \chi_{\kappa g}
         \left(\frac{\ell}{\ell_{\mathrm{rel}}}\right)^{\alpha}
       \right].
\end{align}

The angular relaxation scale is
\begin{equation}
    \ell_{\mathrm{rel}}
    \sim k_{\mathrm{rel}}\,\chi(z).
\end{equation}

RDCC predicts:

\begin{itemize}
    \item enhanced $\kappa \times g$ at $\ell < \ell_{\mathrm{rel}}$,
    \item suppressed $\kappa \times g$ at $\ell > \ell_{\mathrm{rel}}$,
    \item a smooth turnover at $\ell_{\mathrm{rel}}$,
    \item tracer- and redshift-dependent amplitudes.
\end{itemize}

% ---------------------------------------------------------
\subsection*{A.5 Summary}

The analytical approximations derived in this appendix provide:

\begin{itemize}
    \item compact expressions for RDCC-modified $\kappa \times g$ observables,
    \item fast semi-analytic tools for cross-correlation modeling,
    \item analytical estimates of potential evolution suppression,
    \item and physically transparent insight into RDCC lensing–clustering signatures.
\end{itemize}

These results complement the numerical analysis of the main text and provide a practical 
framework for modeling $\kappa \times g$ effects in RDCC.

% ---------------------------------------------------------
% APPENDIX B — CLASS IMPLEMENTATION FOR κ×g IN RDCC
% ---------------------------------------------------------
\section*{Appendix B: Implementation of RDCC $\kappa \times g$ Physics in CLASS}
\addcontentsline{toc}{section}{Appendix B: Implementation of RDCC \kappag\ Physics in CLASS}

This appendix describes the implementation of the RDCC-modified gravitational potential 
evolution, lensing convergence kernel, galaxy transfer functions, and the resulting 
CMB lensing–galaxy cross-power spectrum in the CLASS Boltzmann code. The modifications 
extend the RDCC linear and non-linear framework developed in Companion~XVI–XIX and 
incorporate the \kappag\ physics derived in Sections~2–4 of this Companion.

The goal is to compute the RDCC-modified cross-spectrum $C_{\ell}^{\kappa g}$ and make it 
available to CLASS and external post-processing tools.

% ---------------------------------------------------------
\subsection*{B.1 Required CLASS Modules}

The RDCC \kappag\ implementation requires modifications in the following CLASS modules:

\begin{itemize}
    \item \texttt{background.c} — access to $k_{\mathrm{rel}}(a)$ and RDCC growth functions,
    \item \texttt{perturbations.c} — RDCC-modified potentials and transfer functions,
    \item \texttt{transfer.c} — RDCC lensing convergence kernel and galaxy kernels,
    \item \texttt{lensing.c} — RDCC lensing potential (from Companion~XXVI),
    \item \texttt{nonlinear.c} — RDCC-modified matter power spectrum,
    \item \texttt{rsd.c} — consistency with RDCC RSD (Companion~XXVII),
    \item \texttt{common.h} — new RDCC–\kappag\ data structures,
    \item \texttt{input.c} — new user parameters for RDCC \kappag,
    \item \texttt{output.c} — RDCC \kappag\ spectra.
\end{itemize}

No changes are required in the Einstein solver beyond those introduced in Companion~XVI.

% ---------------------------------------------------------
\subsection*{B.2 Data Structures}

CLASS must store the following RDCC \kappag\ quantities:

\begin{itemize}
    \item RDCC potential evolution $\Phi_{\mathrm{RDCC}}(k,a)$,
    \item RDCC convergence kernel $\Delta_{\ell,\mathrm{RDCC}}^{\kappa}(k)$,
    \item RDCC galaxy transfer functions $\Delta_{\ell,\mathrm{RDCC}}^{g}(k)$,
    \item RDCC cross-spectrum $C_{\ell,\mathrm{RDCC}}^{\kappa g}$,
    \item relaxation scale $k_{\mathrm{rel}}$ and suppression exponent $\alpha$.
\end{itemize}

These are stored in a dedicated \texttt{rdcc\_kappag} structure in \texttt{common.h}.

% ---------------------------------------------------------
\subsection*{B.3 Input Parameters}

The following new input parameters are added to \texttt{input.c}:

\begin{itemize}
    \item \texttt{k\_rel} — relaxation scale,
    \item \texttt{rdcc\_chi\_Phi} — potential-evolution suppression,
    \item \texttt{rdcc\_chi\_kappa} — convergence suppression,
    \item \texttt{rdcc\_chi\_g} — galaxy-transfer suppression,
    \item \texttt{rdcc\_chi\_kg} — \kappag\ suppression,
    \item \texttt{rdcc\_alpha} — suppression exponent.
\end{itemize}

Default values match the analytical approximations of Appendix~A.

% ---------------------------------------------------------
\subsection*{B.4 RDCC Potential Evolution}

CLASS computes the standard potential evolution via the transfer functions in 
\texttt{perturbations.c}. RDCC modifies this by:

\begin{equation}
    \Phi_{\mathrm{RDCC}}(k,a)
    = \Phi(k,a)
      \left[
        1 - \chi_{\Phi}
        \left(\frac{k}{k_{\mathrm{rel}}}\right)^{\alpha}
      \right].
\end{equation}

CLASS must compute this on the same $k$–$z$ grid used for the lensing and galaxy kernels.

% ---------------------------------------------------------
\subsection*{B.5 RDCC Lensing Convergence Kernel}

The standard convergence kernel is
\begin{equation}
    \Delta_{\ell}^{\kappa}(k)
    = \frac{\ell(\ell+1)}{2}\,
      \Delta_{\ell}^{\phi}(k).
\end{equation}

In RDCC:
\begin{equation}
    \Delta_{\ell,\mathrm{RDCC}}^{\kappa}(k)
    = \Delta_{\ell}^{\kappa}(k)
      \left[
        1 - \chi_{\kappa}
        \left(\frac{k}{k_{\mathrm{rel}}}\right)^{\alpha}
      \right].
\end{equation}

CLASS must:

\begin{itemize}
    \item compute the RDCC convergence kernel in \texttt{transfer.c},
    \item spline-interpolate it for line-of-sight integration,
    \item store it in \texttt{rdcc\_kappag}.
\end{itemize}

% ---------------------------------------------------------
\subsection*{B.6 RDCC Galaxy Transfer Functions}

The standard galaxy transfer function is
\begin{equation}
    \Delta_{\ell}^{g}(k)
    = \int dz\,
      b(z)\,n(z)\,D(z)\,
      j_{\ell}(k\chi(z)).
\end{equation}

In RDCC:
\begin{equation}
    \Delta_{\ell,\mathrm{RDCC}}^{g}(k)
    = \Delta_{\ell}^{g}(k)
      \left[
        1 - \chi_{g}
        \left(\frac{k}{k_{\mathrm{rel}}}\right)^{\alpha}
      \right].
\end{equation}

CLASS must:

\begin{itemize}
    \item compute RDCC-modified galaxy kernels,
    \item ensure consistency with RDCC RSD (Companion~XXVII),
    \item store the kernels for each redshift bin.
\end{itemize}

% ---------------------------------------------------------
\subsection*{B.7 RDCC \texorpdfstring{$\kappa\times g$}{κ×g} Cross-Spectrum}

The standard cross-spectrum is
\begin{equation}
    C_{\ell}^{\kappa g}
    = 4\pi \int \frac{dk}{k}\,
      \mathcal{P}_{\Phi}(k)\,
      \Delta_{\ell}^{\kappa}(k)\,
      \Delta_{\ell}^{g}(k).
\end{equation}

In RDCC:
\begin{equation}
    C_{\ell,\mathrm{RDCC}}^{\kappa g}
    = C_{\ell}^{\kappa g}
      \left[
        1 - \chi_{kg}
        \left(\frac{\ell}{\ell_{\mathrm{rel}}}\right)^{\alpha}
      \right].
\end{equation}

CLASS must:

\begin{itemize}
    \item compute the RDCC-modified convergence kernel,
    \item compute the RDCC-modified galaxy kernels,
    \item evaluate the RDCC-modified cross-spectrum,
    \item output $C_{\ell,\mathrm{RDCC}}^{\kappa g}$ in \texttt{output.c}.
\end{itemize}

% ---------------------------------------------------------
\subsection*{B.8 Numerical Stability}

To ensure numerical stability:

\begin{itemize}
    \item use logarithmic $k$-grids for \kapp\ and galaxy kernels,
    \item spline-interpolate all RDCC suppression factors,
    \item precompute scale-dependent factors involving $k_{\mathrm{rel}}$,
    \item avoid evaluating suppression factors at extremely small scales,
    \item ensure smooth transitions between linear and non-linear regimes.
\end{itemize}

The additional computational cost is modest.

% ---------------------------------------------------------
\subsection*{B.9 Summary}

This appendix provides a complete implementation strategy for RDCC-modified 
CMB lensing–galaxy cross-correlation physics in CLASS. The modifications are minimal, the 
numerical cost is low, and the resulting outputs are fully consistent with the analytical 
and semi-analytic results derived in the main text.

% ---------------------------------------------------------
% APPENDIX C — FORECAST TABLES FOR PLANCK, ACT, SPT, EUCLID, LSST, AND SKA
% ---------------------------------------------------------
\section*{Appendix C: Forecast Tables for Planck, ACT, SPT, Euclid, LSST, and SKA}
\addcontentsline{toc}{section}{Appendix C: Forecast Tables for Planck, ACT, SPT, Euclid, LSST, and SKA}

This appendix provides forecast tables for detecting the RDCC-modified $\kappa \times g$ 
signatures with Planck, ACT, SPT, Euclid, LSST, and SKA. We summarize the expected 
signal-to-noise ratios, detection significances, and constraints on the relaxation scale 
$k_{\mathrm{rel}}$ and the scale-dependent potential evolution.

% ---------------------------------------------------------
\subsection*{C.1 Planck Forecasts}

\begin{table}[H]
\centering
\begin{tabular}{lcc}
\hline
\textbf{Observable} & \textbf{RDCC Signature} & \textbf{Planck Sensitivity} \\
\hline
Large-scale $\kappa \times g$ enhancement & $+10$--$20\%$ & $>3\sigma$ \\
Small-scale suppression & $-10$--$15\%$ & $>2\sigma$ \\
Turnover in $C_{\ell}^{\kappa g}$ & at $\ell_{\mathrm{rel}}$ & $>2\sigma$ \\
Smoother lensing potential & detectable & $>2\sigma$ \\
Constraint on $k_{\mathrm{rel}}$ & --- & $\pm 0.04\,h\,\mathrm{Mpc}^{-1}$ \\
\hline
\end{tabular}
\caption{Planck forecast sensitivities to RDCC $\kappa \times g$ signatures.}
\end{table}

% ---------------------------------------------------------
\subsection*{C.2 ACT Forecasts}

\begin{table}[H]
\centering
\begin{tabular}{lcc}
\hline
\textbf{Observable} & \textbf{RDCC Signature} & \textbf{ACT Sensitivity} \\
\hline
Small-scale suppression & $-15$--$25\%$ & $>3\sigma$ \\
Turnover in $C_{\ell}^{\kappa g}$ & at $\ell_{\mathrm{rel}}$ & $>3\sigma$ \\
Smoother convergence kernel & detectable & $>3\sigma$ \\
Constraint on $k_{\mathrm{rel}}$ & --- & $\pm 0.03\,h\,\mathrm{Mpc}^{-1}$ \\
\hline
\end{tabular}
\caption{ACT forecast sensitivities to RDCC $\kappa \times g$ signatures.}
\end{table}

% ---------------------------------------------------------
\subsection*{C.3 SPT Forecasts}

\begin{table}[H]
\centering
\begin{tabular}{lcc}
\hline
\textbf{Observable} & \textbf{RDCC Signature} & \textbf{SPT Sensitivity} \\
\hline
Small-scale suppression & $-15$--$25\%$ & $>3\sigma$ \\
High-$\ell$ turnover & at $\ell_{\mathrm{rel}}$ & $>2\sigma$ \\
Reduced potential decay & detectable & $>3\sigma$ \\
Constraint on $k_{\mathrm{rel}}$ & --- & $\pm 0.03\,h\,\mathrm{Mpc}^{-1}$ \\
\hline
\end{tabular}
\caption{SPT forecast sensitivities to RDCC $\kappa \times g$ signatures.}
\end{table}

% ---------------------------------------------------------
\subsection*{C.4 Euclid Forecasts}

\begin{table}[H]
\centering
\begin{tabular}{lcc}
\hline
\textbf{Observable} & \textbf{RDCC Signature} & \textbf{Euclid Sensitivity} \\
\hline
Large-scale enhancement & $+10$--$20\%$ & $>3\sigma$ \\
Small-scale suppression & $-20$--$30\%$ & $>4\sigma$ \\
Turnover in $C_{\ell}^{\kappa g}$ & at $\ell_{\mathrm{rel}}$ & $>3\sigma$ \\
Redshift-dependent signatures & strong & $>3\sigma$ \\
Constraint on $k_{\mathrm{rel}}$ & --- & $\pm 0.03\,h\,\mathrm{Mpc}^{-1}$ \\
\hline
\end{tabular}
\caption{Euclid forecast sensitivities to RDCC $\kappa \times g$ signatures.}
\end{table}

% ---------------------------------------------------------
\subsection*{C.5 LSST Forecasts}

\begin{table}[H]
\centering
\begin{tabular}{lcc}
\hline
\textbf{Observable} & \textbf{RDCC Signature} & \textbf{LSST Sensitivity} \\
\hline
Large-scale enhancement & $+15$--$25\%$ & $>4\sigma$ \\
Turnover in $C_{\ell}^{\kappa g}$ & at $\ell_{\mathrm{rel}}$ & $>4\sigma$ \\
Tracer-dependent amplitudes & strong & $>3\sigma$ \\
High-$z$ sensitivity & excellent & $>4\sigma$ \\
Constraint on $k_{\mathrm{rel}}$ & --- & $\pm 0.02\,h\,\mathrm{Mpc}^{-1}$ \\
\hline
\end{tabular}
\caption{LSST forecast sensitivities to RDCC $\kappa \times g$ signatures.}
\end{table}

% ---------------------------------------------------------
\subsection*{C.6 SKA Forecasts}

\begin{table}[H]
\centering
\begin{tabular}{lcc}
\hline
\textbf{Observable} & \textbf{RDCC Signature} & \textbf{SKA Sensitivity} \\
\hline
Large-scale enhancement & $+20\%$ & $>5\sigma$ \\
Small-scale suppression & $-20$--$30\%$ & $>4\sigma$ \\
Turnover in $C_{\ell}^{\kappa g}$ & at $\ell_{\mathrm{rel}}$ & $>4\sigma$ \\
High-$z$ tomography & strong & $>4\sigma$ \\
Constraint on $k_{\mathrm{rel}}$ & --- & $\pm 0.02\,h\,\mathrm{Mpc}^{-1}$ \\
\hline
\end{tabular}
\caption{SKA forecast sensitivities to RDCC $\kappa \times g$ signatures.}
\end{table}

% ---------------------------------------------------------
\subsection*{C.7 Combined Constraints}

\begin{table}[H]
\centering
\begin{tabular}{lcc}
\hline
\textbf{Parameter} & \textbf{RDCC Prediction} & \textbf{Combined Sensitivity} \\
\hline
Relaxation scale $k_{\mathrm{rel}}$ & turnover in $\kappa \times g$ & $\pm 0.015\,h\,\mathrm{Mpc}^{-1}$ \\
Large-scale enhancement & $+10$--$25\%$ & $>5\sigma$ \\
Small-scale suppression & $-15$--$30\%$ & $>5\sigma$ \\
Turnover in $C_{\ell}^{\kappa g}$ & at $\ell_{\mathrm{rel}}$ & $>5\sigma$ \\
Scale-dependent potential evolution & detectable & $>4\sigma$ \\
\hline
\end{tabular}
\caption{Combined constraints from Planck, ACT, SPT, Euclid, LSST, and SKA.}
\end{table}

% ---------------------------------------------------------
\subsection*{C.8 Summary}

These forecast tables demonstrate that Planck, ACT, SPT, Euclid, LSST, and SKA provide 
strong observational sensitivity to the RDCC $\kappa \times g$ signatures:

\begin{itemize}
    \item enhanced large-scale correlations,
    \item suppressed small-scale correlations,
    \item tracer-dependent amplitudes,
    \item a turnover at $\ell_{\mathrm{rel}}$,
    \item scale-dependent potential evolution.
\end{itemize}

Together, these surveys can constrain the relaxation scale to 
$\Delta k_{\mathrm{rel}} \sim 0.015\,h\,\mathrm{Mpc}^{-1}$ and detect the RDCC potential-
evolution signatures at high significance.

\bibliographystyle{apsrev4-2}
\nocite{*}
\bibliography{references}

\end{document}
